% Stata exports tabular-only bodies (no \begin{table}); this wrapper owns float, caption, label, notes, width.
\providecommand{\ChapterTableNoteText}[1]{\parbox[t]{0.94\textwidth}{#1}}
\providecommand{\InputStataRegTable}[4]{%
\begin{table}[htbp]
\centering
\caption{#1}%
\label{#2}%
\begin{threeparttable}
\small
\resizebox{\textwidth}{!}{%
\input{#3}%
}\n\begin{tablenotes}[flushleft]
\footnotesize
\item \ChapterTableNoteText{#4}%
\end{tablenotes}
\end{threeparttable}
\end{table}%
}
% Tighter layout for wide panel tables.
\providecommand{\InputStataRegTableCompact}[4]{%
\begin{table}[htbp]
\centering
\caption{#1}%
\label{#2}%
\begin{threeparttable}
\begingroup
\footnotesize
\setlength{\tabcolsep}{4pt}%
\renewcommand{\arraystretch}{0.94}%
% Use adjustbox with an explicit height limit for compact tables to avoid overflowing appendix pages.
\begin{adjustbox}{max totalsize={\textwidth}{0.65\textheight}}
\input{#3}%
\end{adjustbox}
\endgroup
\begin{tablenotes}[flushleft]
\footnotesize
\item \ChapterTableNoteText{#4}%
\end{tablenotes}
\end{threeparttable}
\end{table}%
}

%%%%%%%%%%%%%%%%%%%%%%%%%%%%%%%%%%%%%%%%%%%%%%%%%%%%%%%%%%%%%%%%%%%%%%%%%%%
\section{Determinants of trust}
\label{sec:appendix}

\par The focus of this paper is the general trust measure. Table~\ref{tab:42_trust_determinants} describes the set of controls and determinants of general trust. The first column includes standard[...] 

% Replaced the single macro call with the compact wrapper (height-limited) to prevent overflow on the appendix page.
\InputStataRegTableCompact{\textbf{Determinants of General Trust}}{tab:42_trust_determinants}{../../Code/Analysis/Tables/42_table1_trust_determinants.tex}{OLS. Heteroskedasticity-robust standard errors. Column~\textit{(1)}: demographics, depression, and labor force. Column~\textit{(2)}: adds population size, wealth deciles, region, born in U.S., and six \say{Potential for Deception} items. Column~\textit{(3)}: adds health conditions, Medicare/Medicaid, life insurance, and marital history. Column~\textit{(4)}: same as column~\textit{(3)} plus an indicator for any bequest. Age-bin fixed effects are included in all columns but omitted from the body of the table.}
\FloatBarrier

\par As \cite{Alesina2000} suggests, demographic variables tend to determine trust. In the HRS data, the estimated coefficients for non-Hispanic Black and for marital status have the strongest sig[...]

%%%%%%%%%%%%%%%%%%%%%%%%%%%%%%%%%%%%%%%%%%%%%%%%%%%%%%%%%%%%%%%%%%%%%%%%%%%
\section{Race coefficients without trust in panel models}

\par Table~\ref{tab:47_race_no_trust_panel} extends the no-trust race comparison to the panel setting. The FE columns report second-stage cross-sectional regressions of the estimated household fix[...]

\InputStataRegTableCompact{\textbf{Race Coefficients in No-Trust FE and CRE Specifications}}{tab:47_race_no_trust_panel}{../../Code/Analysis/Tables/47_table2_race_no_trust_panel.tex}{FE columns report second-stage OLS of the estimated household fixed effect from no-trust FE specifications; CRE columns report no-trust CRE regressions of 5\% winsorized returns to net wealth. In each model family, the reduced column includes race with the baseline return controls, while the full column adds the extended demographic controls. Only race coefficients are displayed.}
\FloatBarrier

%%%%%%%%%%%%%%%%%%%%%%%%%%%%%%%%%%%%%%%%%%%%%%%%%%%%%%%%%%%%%%%%%%%%%%%%%%%
\section{Measures of income}

\par Labor income is a narrow measure only capturing earnings and unemployment income. Total income is a more broad measure including retirement and capital income. Figure~\ref{fig:income_over_time} shows that, for the time period, mean incomes are relatively flat over the period. Moreover, the discrepancy between mean labor and total income is likely driven by the oversampling of older and retired households by the HRS.

\begin{figure}[H]
\centering
\includegraphics[width=0.8\textwidth]{../../Code/Descriptive/Figures/income_over_time.png}
\caption{Income over time}
\label{fig:income_over_time}
\end{figure}
\FloatBarrier

\par In particular, I pooled the observations of the survey together to get a sense of the distribution of incomes measured in the survey. This is captured by Table~\ref{tab:42_income_tabstat}.

\InputStataRegTableCompact{\textbf{Income Distribution Summary Statistics}}{tab:42_income_tabstat}{../../Code/Analysis/Tables/42_table2_income_tabstat.tex}{Real USD, winsorized; person-years restricted to the same pooled regression sample used for return descriptives in the main text (Tables~\ref{tab:39_r5_age_bin} and~\ref{tab:39_r5_educ_cat}).}
\FloatBarrier

\par An important finding in the literature on earnings in the U.S. is that it is generally hump-shaped over the lifecycle. It is a good sign then that Figure~\ref{fig:income_by_agegroup_2020} displays this pattern.

\begin{figure}[H]
\centering
\includegraphics[width=0.8\textwidth]{../../Code/Descriptive/Figures/income_by_agegroup_2020.png}
\caption{Income by age group (2020)}
\label{fig:income_by_agegroup_2020}
\end{figure}
\FloatBarrier

\par Another empirical finding in the literature on earnings is that on average individuals with more education earn higher income. The trend captured in Table~\ref{tab:42_income_by_educ} suggests that the earnings data is in line with what one would expect regarding earnings for a representative survey in the U.S.\footnote{Although the HRS oversamples older households, I use the provided respondent-level weights for this interpretation of the summary statistics of the data.}

\InputStataRegTableCompact{\textbf{Mean Income by Education Group}}{tab:42_income_by_educ}{../../Code/Analysis/Tables/42_table3_income_by_educ.tex}{Real USD, winsorized; same person-year sample as Table~\ref{tab:42_income_tabstat}. No hs $=$ $<$12y; Hs $=$ 12y; Some college $=$ 13--15y; 4yr $=$ 16y; Grad $=$ 17+y.}
